\chapter{What an Algorithm Really Is}

You already know more algorithms than you think. When you sort a hand of playing
cards, you are running one. When you look up a word in a dictionary by flipping to
the middle and then deciding ``earlier'' or ``later,'' you are running another --
and a surprisingly good one. When you give a friend turn-by-turn directions to your
house, you are writing one down. The word \emph{algorithm} sounds technical and
cold, but the thing it names is something you have been doing since childhood: a
finite recipe of unambiguous steps that turns a question into an answer.

This chapter is about taking that everyday notion and sharpening it until it is
precise enough to reason about. We are going to ask, slowly and carefully: what
exactly \emph{is} an algorithm? What does it mean for one to be \emph{correct}?
Why do we care how \emph{fast} it is, when computers keep getting faster anyway?
And how do we write algorithms down so that another human -- or our future self --
can understand them without drowning in the fiddly details of any one programming
language?

We will not prove deep theorems yet. The goal is to build the vocabulary and the
habits of mind that the rest of the book leans on. By the end you will be able to
read a piece of pseudocode, say precisely what problem it solves, argue that it
solves that problem, and estimate how its running time grows as the input grows.
Those four skills -- read, specify, justify, estimate -- are the spine of
everything that follows.

\chapterline

\section{Recipes, and why precision matters}

Start with the friendliest analogy and then watch it break, because the break is
where the real ideas live.

\begin{intuitionbox}
An algorithm is a recipe. It takes some \emph{ingredients} (the input), follows a
fixed sequence of \emph{steps}, and produces a \emph{dish} (the output). The two
demands we place on a recipe are the two demands we place on an algorithm:
every step must be something you can actually do, and the whole thing must
eventually finish.
\end{intuitionbox}

Suppose I hand you this ``recipe'' for finding the largest number in a list:

\begin{quote}
\emph{Look at all the numbers and write down the biggest one.}
\end{quote}

Is that an algorithm? It has an input (a list of numbers) and an output (the biggest
one), and a human can ``run'' it. But look closer. ``Look at all the numbers'' --
in what order? All at once? A computer cannot glance at a thousand numbers
simultaneously the way your eye sweeps a page. ``Write down the biggest one'' --
how did you decide which is biggest? The instruction quietly assumes you already
know how to do the very thing we are trying to describe. It is a wish, not a recipe.

Compare it with this:

\begin{quote}
\emph{Set a variable \texttt{best} to the first number. Then, going left to right
through the rest of the numbers one at a time, whenever you meet a number larger
than \texttt{best}, replace \texttt{best} with that number. When you reach the end,
\texttt{best} holds the answer.}
\end{quote}

This version never asks you to do anything you do not already know how to do:
read a number, compare two numbers, copy a value. There is no hidden magic step.
And it tells you the order. That is the difference between a vague instruction and
an algorithm. The whole discipline of this subject is, in a sense, the relentless
removal of hidden magic steps.

\begin{focusbox}
An algorithm is a \textbf{finite}, \textbf{unambiguous}, \textbf{effective}
sequence of steps that transforms an input into the desired output.
\emph{Finite}: it stops. \emph{Unambiguous}: at every moment it is clear exactly
what to do next. \emph{Effective}: each step is a genuinely doable operation, not a
restatement of the problem.
\end{focusbox}

\begin{historybox}
The word \emph{algorithm} comes from the name of the 9th-century Persian
mathematician Muhammad ibn Musa al-Khwarizmi, whose book on calculation gave Europe
its methods for arithmetic with Hindu-Arabic numerals. (His other great book,
\emph{al-jabr}, gave us the word \emph{algebra}.) The modern, mathematically precise
notion of ``a step a machine can carry out'' was pinned down in the 1930s by Alan
Turing and Alonzo Church. This chapter parallels \textbf{CLRS Chapter~1, ``The Role
of Algorithms in Computing''}; the deeper theory of what is and is not computable is
hinted at there and developed elsewhere.
\end{historybox}

\section{Problems versus instances}

Here is a distinction that beginners blur and experts never do, so let us nail it
down early. There is a difference between a \emph{problem} and an \emph{instance} of
that problem.

A \textbf{problem} is the general task, described in terms of \emph{any} possible
input. An \textbf{instance} is one specific input you might be handed.

\begin{definition}[ch01:problem-instance]
A \emph{computational problem} is specified by describing (i) the set of allowed
\emph{inputs}, often called \emph{instances}, and (ii) for each input, what counts
as a correct \emph{output}. An \emph{instance} of the problem is a single concrete
input drawn from that set.
\end{definition}

Take the sorting problem. As a \emph{problem} it reads:

\begin{formalbox}
\textbf{The Sorting Problem.}\\
\textbf{Input:} a sequence of $n$ numbers $\langle a_1, a_2, \dots, a_n\rangle$.\\
\textbf{Output:} a reordering $\langle a'_1, a'_2, \dots, a'_n\rangle$ of the input
sequence such that $a'_1 \le a'_2 \le \cdots \le a'_n$.
\end{formalbox}

That definition does not mention any particular list. It speaks of \emph{every}
list at once. An \emph{instance} of the sorting problem is something concrete, like
$\langle 31, 41, 59, 26, 41, 58 \rangle$, whose correct output is the single
sequence $\langle 26, 31, 41, 41, 58, 59 \rangle$.

\begin{intuitionbox}
Think of a problem as a giant lookup table you can never finish writing: down the
left column is every conceivable input, and in the right column is the correct
answer for that input. An \emph{instance} is one row. An \emph{algorithm} is a
finite machine that, given any row's left entry, computes its right entry -- without
ever storing the impossible infinite table.
\end{intuitionbox}

Figure~\ref{ch01:problem-table} draws that mental picture for the sorting problem.

\begin{figure}[H]
\centering
\begin{tikzpicture}[
    node distance=0pt,
    row/.style={draw,minimum height=0.62cm,inner xsep=8pt,font=\footnotesize,anchor=west},
    lcol/.style={row,minimum width=4.2cm,fill=sourceblue!8},
    rcol/.style={row,minimum width=4.2cm,fill=targetgreen!10}]
  % header
  \node[lcol,fill=sourceblue!20,font=\footnotesize\bfseries] (h1) at (0,0)
        {input (instance)};
  \node[rcol,fill=targetgreen!22,font=\footnotesize\bfseries] (h2) at (4.2,0)
        {correct output};
  % rows
  \node[lcol] (l1) at (0,-0.62) {$\langle 3,1,2\rangle$};
  \node[rcol] (r1) at (4.2,-0.62) {$\langle 1,2,3\rangle$};
  \node[lcol] (l2) at (0,-1.24) {$\langle 31,41,59,26\rangle$};
  \node[rcol] (r2) at (4.2,-1.24) {$\langle 26,31,41,59\rangle$};
  \node[lcol] (l3) at (0,-1.86) {$\langle 5\rangle$};
  \node[rcol] (r3) at (4.2,-1.86) {$\langle 5\rangle$};
  \node[lcol,fill=quotientorange!18,draw=quotientorange,line width=1pt] (l4) at (0,-2.48)
        {$\langle 31,41,59,26,41,58\rangle$};
  \node[rcol,fill=quotientorange!18,draw=quotientorange,line width=1pt] (r4) at (4.2,-2.48)
        {$\langle 26,31,41,41,58,59\rangle$};
  \node[lcol] (l5) at (0,-3.10) {$\vdots$};
  \node[rcol] (r5) at (4.2,-3.10) {$\vdots$};
  % the algorithm arrow
  \draw[->,thick,sourceblue] (l4.east)+(0,0) -- (r4.west)
        node[midway,above=1pt,font=\tiny\bfseries,sourceblue] {algorithm};
  \node[font=\footnotesize,quotientorange,anchor=west] at (8.7,-2.48)
        {$\leftarrow$ one instance (one row)};
\end{tikzpicture}
\caption{The sorting problem as an infinite lookup table: every possible input on the
left, its unique correct output on the right. The highlighted row is one
\emph{instance}. An \emph{algorithm} is a finite machine that computes the right entry
from the left entry for \emph{any} row, without ever storing the table.}
\label{ch01:problem-table}
\end{figure}

Why does this distinction earn its own section? Because almost every confusion later
in the book dissolves once you ask, ``Wait, am I talking about the problem or about
one instance?'' When we say an algorithm is \emph{correct}, we mean it works on
\emph{every} instance, not just the one we tried. When we say an algorithm is
\emph{slow}, we usually mean there exists some family of instances on which it bogs
down. And when we compare two algorithms, we compare them as solvers of the same
\emph{problem}, even though they may behave very differently on a given
\emph{instance}.

\begin{trapbox}
``I ran it on $\langle 3, 1, 2 \rangle$ and it gave $\langle 1, 2, 3 \rangle$, so my
sorting algorithm is correct.'' No -- that shows it is correct on \emph{one
instance}. Correctness is a claim about \emph{all} instances. A single successful
example can never establish it; a single failing example can demolish it. Testing
finds bugs; it does not prove their absence.
\end{trapbox}

\section{What ``correct'' actually means}

We have been using the word \emph{correct} freely. Let us make it sharp.

\begin{definition}[ch01:correctness]
An algorithm is \emph{correct} for a problem if, for \emph{every} instance of that
problem, it halts and returns a correct output for that instance. An algorithm that
halts with a correct answer on some inputs but on others either halts with a wrong
answer or never halts at all is \emph{incorrect}.
\end{definition}

Two words in there deserve a spotlight: \emph{every} and \emph{halts}.

\emph{Every} we have already met -- correctness quantifies over all instances. The
new word is \emph{halts}. An algorithm that runs forever has not solved your
problem, no matter how clever it looks, because you never get an answer. So
correctness bundles two promises together: the algorithm \emph{terminates}, and when
it terminates the answer is \emph{right}. We will return to both promises again and
again, usually proving them with a tool called a \emph{loop invariant}, which we
preview at the end of this chapter and use in earnest from the next one onward.

\begin{remark}[ch01:approximation]
For some hard problems we deliberately relax ``returns a correct output'' to
``returns an output provably close to correct.'' Those are \emph{approximation
algorithms}, a whole subject of their own. And some algorithms flip a coin
internally and are only correct \emph{with high probability}; those are
\emph{randomized algorithms}. Both are legitimate and useful -- but in each case the
relaxation is stated openly and precisely. Vagueness is never the relaxation.
\end{remark}

\begin{aestheticbox}
Why be so fussy about ``every instance''? Because software runs on inputs nobody
imagined. A date library that is correct on every day you tested but loops forever
on the 29th of February in a leap year is not ``almost correct'' -- it is the bug
that takes down an airline at 2~a.m. The habit of asking ``for which inputs might
this fail?'' is the single most valuable reflex this book can give you. Correctness
is not academic pedantry; it is the difference between code that works and code that
works \emph{until it doesn't}.
\end{aestheticbox}

\section{A first algorithm, written out properly}

Now we turn the ``find the maximum'' recipe from earlier into something we can
analyze. Along the way we introduce the way this book writes algorithms down --
\emph{pseudocode} -- and explain the conventions as we go.

\begin{algorithm}[H]
\caption{\textsc{Find-Max}: largest element of a list}
\KwIn{An array $A[1\,..\,n]$ of $n$ numbers, with $n \ge 1$.}
\KwOut{The value of the largest element of $A$.}
$best \gets A[1]$\;
\For{$i \gets 2$ \KwTo $n$}{
  \If{$A[i] > best$}{
    $best \gets A[i]$\;
  }
}
\Return $best$\;
\end{algorithm}

Read it slowly. The first line copies the first element into a variable named
$best$, our running champion. The loop then walks the index $i$ from $2$ up to $n$,
inspecting each remaining element exactly once. Whenever the current element beats
the champion, it becomes the new champion. After the loop, $best$ holds the answer,
and we return it.

\subsection{Pseudocode conventions}

Pseudocode is not a programming language. It is a controlled form of English-plus-math
whose only goal is to be unambiguous to a human reader. We strip away everything
that does not matter to the \emph{algorithm} -- semicolons, type declarations,
memory management -- and keep everything that does. Here are the conventions this
book uses; they match CLRS closely.

\begin{itemize}
  \item \textbf{Assignment} is written $x \gets v$, read ``$x$ gets $v$.'' It
        copies the value $v$ into the variable $x$. We avoid the symbol $=$ for
        assignment so that $=$ can mean plain mathematical equality.
  \item \textbf{Arrays} are written $A[1\,..\,n]$ and indexed from $1$, so $A[1]$ is
        the first element and $A[n]$ the last. (Many real languages index from
        $0$; the ideas are identical, the arithmetic shifts by one.) The notation
        $A[i\,..\,j]$ means the slice from index $i$ through $j$ inclusive.
  \item \textbf{Blocks} are shown by indentation, not braces. The body of the
        \textbf{for} loop is everything indented beneath it.
  \item \textbf{Loops} come in \textbf{for} (a known number of repetitions) and
        \textbf{while} (repeat as long as a condition holds) flavors. After a
        \textbf{for} loop ends, the loop variable is considered to hold the value
        that first failed the test -- a small convention worth remembering.
  \item \textbf{Conditionals} are \textbf{if} / \textbf{else}. \textbf{Return}
        immediately stops the algorithm and hands back a value.
  \item We freely use ordinary math -- $\le$, $\lfloor n/2 \rfloor$, ``swap $x$ and
        $y$'' -- whenever it is clearer than spelling out machine steps.
\end{itemize}

\begin{focusbox}
The point of pseudocode is communication, not execution. If a human reader could
follow your steps and get the right answer by hand, your pseudocode is good --
even if no compiler would accept it. Aim for the clarity of a well-written proof,
not the syntax of a particular language.
\end{focusbox}

\subsection{Tracing it on an instance}

Abstract steps come alive when you run them by hand. Trace \textsc{Find-Max}
on the instance $A = \langle 4, 9, 3, 9, 7 \rangle$, with $n = 5$. The table records
the state \emph{after} each iteration of the loop.

\begin{center}
\begin{tabular}{c c c c}
\toprule
Step & $i$ & $A[i]$ & $best$ after this step \\
\midrule
initialize & --  & --        & $4$ \\
loop       & $2$ & $9$       & $9$ \quad (since $9 > 4$) \\
loop       & $3$ & $3$       & $9$ \quad (since $3 \not> 9$) \\
loop       & $4$ & $9$       & $9$ \quad (since $9 \not> 9$) \\
loop       & $5$ & $7$       & $9$ \quad (since $7 \not> 9$) \\
\bottomrule
\end{tabular}
\end{center}

We return $9$, which is indeed the largest element. Notice the tie at $i = 4$: the
test is strictly $>$, so an equal value does not displace the champion. That choice
does not matter for the maximum, but in other algorithms the difference between
$>$ and $\ge$ decides whether a sort is \emph{stable} or whether a loop terminates,
so train yourself to notice it.

\begin{labbox}
Trace \textsc{Find-Max} yourself on $A = \langle 5, 5, 5 \rangle$ and on
$A = \langle 8, 1, 6, 9, 2 \rangle$. For each, write the four-column table above.
Then answer: on an array where the maximum sits at the very first position, how many
times does the body of the \textbf{if} (the assignment $best \gets A[i]$) actually
execute? On an array sorted strictly increasing, how many times? This gap between
``never'' and ``every time'' is your first encounter with \emph{best case} versus
\emph{worst case}.
\end{labbox}

\section{Why correctness needs an argument, not a vibe}

\textsc{Find-Max} is simple enough that it ``obviously'' works. But ``obviously''
is exactly the feeling that hides bugs, so we will practice arguing correctness
properly on an easy example, where the machinery is visible and the stakes are low.

The key idea is a \emph{loop invariant}: a statement that is true every time we are
about to test the loop condition, and that, once the loop ends, hands us the result
we want. Think of it as a promise the loop keeps to itself at the top of every pass.

\begin{focusbox}
\textbf{Loop invariant for \textsc{Find-Max}.} Just before the loop runs with a
given value of $i$, the variable $best$ equals the maximum of the first $i-1$
elements, that is, $best = \max\{A[1], A[2], \dots, A[i-1]\}$.
\end{focusbox}

Figure~\ref{ch01:invariant-regions} pictures the promise. At the top of the pass for
a given $i$, the array is split into a settled prefix $A[1\,..\,i-1]$ -- already
scanned, with its maximum captured in $best$ -- and an unexamined suffix
$A[i\,..\,n]$. The whole job of one pass is to fold $A[i]$ into the settled region
and slide the boundary one step to the right.

\begin{figure}[H]
\centering
\begin{tikzpicture}[
    x=0.95cm,
    cell/.style={draw,minimum width=0.95cm,minimum height=0.7cm,inner sep=0pt}]
  % settled prefix
  \foreach \i in {1,2,3,4} {
    \node[cell,fill=targetgreen!16] at (\i,0) {};
  }
  % current element being folded in
  \node[cell,fill=quotientorange!22,draw=quotientorange,line width=1pt] at (5,0) (cur) {};
  % unseen suffix
  \foreach \i in {6,7,8} {
    \node[cell,fill=gray!8] at (\i,0) {};
  }
  % braces
  \draw[decorate,decoration={brace,amplitude=5pt},thick]
    (0.55,0.5) -- (4.45,0.5)
    node[midway,above=6pt,font=\footnotesize] {settled: $best=\max A[1\,..\,i-1]$};
  \draw[decorate,decoration={brace,amplitude=5pt,mirror},thick,gray]
    (5.55,-0.5) -- (8.45,-0.5)
    node[midway,below=6pt,font=\footnotesize,gray] {not yet examined: $A[i+1\,..\,n]$};
  \node[font=\footnotesize,quotientorange] at (5,-0.95) {$A[i]$};
  \node[font=\footnotesize] at (1,-0.55) {$1$};
  \node[font=\footnotesize] at (4,-0.55) {$i\!-\!1$};
  \node[font=\footnotesize] at (8,-0.55) {$n$};
\end{tikzpicture}
\caption{The loop invariant of \textsc{Find-Max} as a picture. Green: the settled
prefix, whose maximum already lives in $best$. Orange: the current element $A[i]$,
about to be compared against $best$ and absorbed. Gray: the suffix not yet looked at.
Each pass moves the green/orange boundary one cell right; when green covers the whole
array, $best$ is the answer.}
\label{ch01:invariant-regions}
\end{figure}

To use an invariant we check three things -- the same three every time, so they will
soon feel routine.

\begin{itemize}
  \item \textbf{Initialization.} The invariant holds before the first iteration.
        Here, before the loop, $best = A[1]$ and $i = 2$, so $i - 1 = 1$ and indeed
        $best = \max\{A[1]\}$. True.
  \item \textbf{Maintenance.} If the invariant holds before an iteration, it still
        holds before the next one. Suppose $best = \max\{A[1],\dots,A[i-1]\}$ as we
        enter the body. The body compares $A[i]$ to $best$ and keeps the larger.
        So afterwards $best = \max\{A[1],\dots,A[i-1], A[i]\} = \max\{A[1],\dots,A[i]\}$.
        Then $i$ increases by $1$, restoring the invariant's form with the new $i$.
        True.
  \item \textbf{Termination.} When the loop ends, the invariant plus the reason the
        loop stopped give the answer. The loop stops when $i = n + 1$. The invariant
        then says $best = \max\{A[1], \dots, A[n]\}$ -- exactly the maximum of the
        whole array, which is what we return. Done.
\end{itemize}

\begin{formalbox}
\textbf{Claim.} For every array $A[1\,..\,n]$ with $n \ge 1$, \textsc{Find-Max}
halts and returns $\max\{A[1], \dots, A[n]\}$.\\[2pt]
\textbf{Why it halts:} the loop variable $i$ starts at $2$, rises by exactly $1$ each
pass, and the loop runs only while $i \le n$; after at most $n-1$ passes the
condition fails. \\
\textbf{Why it is right:} the three invariant checks above.
\end{formalbox}

Notice the shape of the argument. \emph{Initialization} is like the base case of an
induction; \emph{maintenance} is like the inductive step; \emph{termination} cashes
the invariant in for the final claim. If that smells like proof by induction, your
nose is excellent -- a loop invariant is induction wearing work clothes. We will lean
on this trio relentlessly, and after a few chapters you will run the three checks
almost without thinking.

\begin{walkbox}
On your next walk, pick a kitchen task you do without thinking -- making tea,
folding laundry, washing a stack of dishes -- and try to write it as careful
pseudocode. Where are the loops? What is the invariant (``after washing $k$ dishes,
the first $k$ are clean and stacked'')? Where did you smuggle in a hidden magic
step that assumes a human's judgment? The discomfort you feel is exactly the gap
between a wish and an algorithm.
\end{walkbox}

\section{Efficiency, and why it matters}

Now for the question that animates most of this book. We have an algorithm that is
correct. Why isn't that the end of the story?

Because correct is not the same as \emph{usable}. Two algorithms can both be
flawlessly correct, return identical answers on every instance, and yet one finishes
in a blink while the other would still be running when the sun burns out. The
resource we care about most is \emph{time}: how many basic steps the algorithm
performs before it stops. (We sometimes also care about \emph{space} -- how much
memory it uses -- but time leads the way.)

\begin{aestheticbox}
You might think: computers double in speed every couple of years, so won't slow
algorithms simply become fast ones if we wait? No -- and this is the central
surprise of the subject. Hardware speed multiplies your throughput by a
\emph{constant} factor. A better algorithm can multiply it by a factor that
\emph{grows with the input size}. On a million-element problem, the right algorithm
can beat the wrong one by a factor of fifty thousand on the very same machine. No
amount of waiting for faster chips closes a gap like that. Algorithms are a
technology, and often the cheapest, most decisive technology in the whole system.
\end{aestheticbox}

\subsection{Counting steps, not seconds}

We deliberately do \emph{not} measure efficiency in seconds. Seconds depend on the
machine, the language, the compiler, the weather inside the chip. Instead we count
\emph{basic operations} -- comparisons, additions, assignments, array accesses --
as a function of the input size $n$. This abstracts away the machine and lets us
compare algorithms on their merits.

For \textsc{Find-Max}, count the comparisons $A[i] > best$. The loop runs for
$i = 2, 3, \dots, n$, which is $n - 1$ iterations, doing one comparison each. So
\textsc{Find-Max} uses exactly $n - 1$ comparisons, on \emph{every} instance of size
$n$. The work grows in step with the input -- double the array, roughly double the
work. We call that \emph{linear} growth, and it is about as good as you can hope for
on a problem that must, at minimum, look at every element at least once.

\subsection{The shape of growth is what matters}

Here is the key idea. When we judge efficiency, we do not fuss over whether an
algorithm does $n - 1$ steps or $n + 3$ or $2n$. We care about the \emph{shape} of
the growth -- linear versus quadratic versus exponential -- because for large
inputs the shape swamps every constant.

The table below shows how different growth shapes balloon as $n$ climbs. Imagine a
machine doing a billion ($10^9$) basic steps per second.

\begin{center}
\begin{tabular}{r r r r r}
\toprule
$n$ & $\log_2 n$ & $n$ & $n^2$ & $2^n$ \\
\midrule
$10$       & $\approx 3$  & $10$       & $100$              & $1024$ \\
$100$      & $\approx 7$  & $100$      & $10{,}000$         & $\approx 10^{30}$ \\
$1{,}000$  & $\approx 10$ & $1{,}000$  & $1{,}000{,}000$    & astronomically huge \\
$10^6$     & $\approx 20$ & $10^6$     & $10^{12}$          & beyond comprehension \\
\bottomrule
\end{tabular}
\end{center}

Read the bottom row as elapsed time on our billion-steps-per-second machine. A
$\log_2 n$ algorithm finishes in about $20$ steps -- instant. A linear one does
$10^6$ steps -- a millisecond. A quadratic one does $10^{12}$ steps -- about
$17$ minutes. And a $2^n$ algorithm on $n = 10^6$? The number of steps has more
digits than there are atoms in the observable universe; it will never finish, on any
machine, ever. The shape of the curve, not the constant out front, decides whether
your program returns before lunch or never.

\begin{intuitionbox}
A constant factor is a one-time discount -- ``half off, forever.'' A better growth
shape is a discount that gets \emph{deeper the bigger your order}. For a small input
the discount barely shows; for a large one it is the whole game. That is why, when
we compare algorithms, we look past the constants and concentrate on the shape.
This squinting has a precise name -- \emph{asymptotic analysis}, written with the
$O$, $\Theta$, $\Omega$ notation -- and it gets its own chapter. For now, get a feel
for the difference between linear and quadratic growth.
\end{intuitionbox}

\begin{figure}[H]
\centering
\begin{tikzpicture}[scale=1.0]
  % axes
  \draw[->,thick] (0,0) -- (6.4,0) node[right] {input size $n$};
  \draw[->,thick] (0,0) -- (0,4.6) node[above] {steps};
  % logarithmic-ish curve
  \draw[sourceblue,very thick,domain=0.12:6,samples=100]
        plot (\x, {0.7*ln(\x*1.6+1)});
  \node[sourceblue] at (5.7,1.55) {$\log n$};
  % linear
  \draw[targetgreen,very thick,domain=0:4.3,samples=2]
        plot (\x, {0.85*\x});
  \node[targetgreen] at (4.0,3.9) {$n$};
  % quadratic
  \draw[trapred,very thick,domain=0:2.18,samples=60]
        plot (\x, {0.92*\x*\x});
  \node[trapred] at (2.45,4.25) {$n^2$};
\end{tikzpicture}
\caption{The same axes, three growth shapes. For small $n$ the curves are tangled
and the constants matter; as $n$ grows the curves fan apart and the \emph{shape}
takes over completely. This fanning is why we judge algorithms by growth, not by a
stopwatch on one input.}
\label{ch01:growth-shapes}
\end{figure}

\section{Comparing two approaches to one task}

Talk of growth shapes stays abstract until you watch two correct algorithms solve the
\emph{same} problem with wildly different efficiency. The problem we will use:
search a list for a target value.

\begin{formalbox}
\textbf{The Searching Problem.}\\
\textbf{Input:} an array $A[1\,..\,n]$ and a value $v$.\\
\textbf{Output:} an index $i$ with $A[i] = v$, or the special value
\textsc{Nil} if $v$ does not appear in $A$.
\end{formalbox}

\subsection{Approach one: linear search}

The obvious method: walk through the array from left to right, and stop the moment
you find $v$.

\begin{algorithm}[H]
\caption{\textsc{Linear-Search}}
\KwIn{An array $A[1\,..\,n]$ and a value $v$.}
\KwOut{An index $i$ with $A[i] = v$, or \textsc{Nil}.}
\For{$i \gets 1$ \KwTo $n$}{
  \If{$A[i] = v$}{
    \Return $i$\;
  }
}
\Return \textsc{Nil}\;
\end{algorithm}

This is correct on any array -- sorted or not -- and uses at most $n$ comparisons. If
$v$ is the first element we stop after one comparison (best case); if $v$ is absent
or last we make all $n$ (worst case). Worst-case work is linear in $n$. Nothing wrong
with that. But watch what happens when we are allowed to assume the array is
\emph{already sorted}.

\subsection{Approach two: binary search}

Suppose $A$ is sorted in increasing order. Now we can be cleverer. Look at the
\emph{middle} element. If it equals $v$, done. If it is \emph{larger} than $v$, then
$v$ -- if present at all -- must lie in the left half, so we throw away the entire
right half and repeat. If it is \emph{smaller}, we throw away the left half. Each
comparison cuts the region still worth searching in half.

\begin{algorithm}[H]
\caption{\textsc{Binary-Search} (on a sorted array)}
\KwIn{A sorted array $A[1\,..\,n]$ with $A[1] \le \cdots \le A[n]$, and a value $v$.}
\KwOut{An index $i$ with $A[i] = v$, or \textsc{Nil}.}
$lo \gets 1$\;
$hi \gets n$\;
\While{$lo \le hi$}{
  $mid \gets \lfloor (lo + hi)/2 \rfloor$\;
  \uIf{$A[mid] = v$}{
    \Return $mid$\;
  }
  \uElseIf{$A[mid] < v$}{
    $lo \gets mid + 1$\;
  }
  \Else{
    $hi \gets mid - 1$\;
  }
}
\Return \textsc{Nil}\;
\end{algorithm}

The variables $lo$ and $hi$ fence off the slice $A[lo\,..\,hi]$ that could still
contain $v$. Each pass examines its midpoint and discards half the slice. The loop
ends either by finding $v$ or by squeezing the slice empty ($lo > hi$), which proves
$v$ is absent.

\begin{intuitionbox}
This is the dictionary trick, made exact. To find ``algorithm'' you do not read from
``aardvark'' onward; you flip to the middle, see you have overshot or undershot, and
flip again into the surviving half. Each flip halves what remains. That halving is
the whole secret, and it is enormously powerful.
\end{intuitionbox}

\subsubsection{Tracing binary search}

Let $A = \langle 2, 5, 8, 12, 16, 23, 38, 56, 72, 91 \rangle$ (so $n = 10$, sorted),
and search for $v = 23$. The state is $(lo, hi, mid, A[mid])$ at the top of each pass.

\begin{center}
\begin{tabular}{c c c c c l}
\toprule
Pass & $lo$ & $hi$ & $mid$ & $A[mid]$ & Decision \\
\midrule
$1$ & $1$ & $10$ & $5$ & $16$ & $16 < 23$, so go right: $lo \gets 6$ \\
$2$ & $6$ & $10$ & $8$ & $56$ & $56 > 23$, so go left: $hi \gets 7$ \\
$3$ & $6$ & $7$  & $6$ & $23$ & found it, \textbf{return} $6$ \\
\bottomrule
\end{tabular}
\end{center}

Three comparisons to find $23$ in a ten-element array. Linear search, looking for the
same value, would have made six. The gap looks modest here -- but it widens
explosively, because of how the halving compounds.

Figure~\ref{ch01:binary-halving} shows the same three passes as shrinking slices.
Each shaded band is the live region $A[lo\,..\,hi]$ still worth searching; the boxed
cell is the midpoint just examined. Watch the live region collapse by roughly half
each time until it lands on the target.

\begin{figure}[H]
\centering
\begin{tikzpicture}[
    x=0.62cm,
    cell/.style={draw,minimum width=0.62cm,minimum height=0.58cm,
                 inner sep=0pt,font=\footnotesize},
    live/.style={fill=sourceblue!16},
    probe/.style={draw=trapred,line width=1pt,fill=trapred!12},
    dead/.style={fill=gray!10,text=gray}]
  \def\vals{2,5,8,12,16,23,38,56,72,91}
  % row labels and three passes; columns 1..10
  \foreach \r/\loi/\hii/\midi in {1/1/10/5, 2/6/10/8, 3/6/7/6} {
    \node[font=\footnotesize\bfseries] at (-1.2,{-1.05*\r+0.6}) {Pass \r};
    \foreach \v [count=\c] in \vals {
      \pgfmathtruncatemacro{\isdead}{(\c<\loi || \c>\hii) ? 1 : 0}
      \pgfmathtruncatemacro{\ismid}{(\c==\midi) ? 1 : 0}
      \ifnum\ismid=1
        \node[cell,probe] at (\c,{-1.05*\r+0.6}) {\v};
      \else
        \ifnum\isdead=1
          \node[cell,dead] at (\c,{-1.05*\r+0.6}) {\v};
        \else
          \node[cell,live] at (\c,{-1.05*\r+0.6}) {\v};
        \fi
      \fi
    }
  }
  % index ruler on top
  \foreach \c in {1,...,10} {
    \node[font=\tiny\itshape,gray] at (\c,1.15) {\c};
  }
\end{tikzpicture}
\caption{Binary search for $v = 23$ in a sorted array of $10$ elements. Shaded cells
are the live slice $A[lo\,..\,hi]$; the red-outlined cell is the midpoint probed that
pass; gray cells have been discarded. Pass~1 probes $A[5]=16<23$ and throws away the
left half; pass~2 probes $A[8]=56>23$ and throws away the right; pass~3 probes
$A[6]=23$ and returns. Three probes, not six.}
\label{ch01:binary-halving}
\end{figure}

\subsection{Counting the halvings}

How many passes can binary search make in the worst case? Each pass at least halves
the size of the live slice. Start with $n$ candidates; after one pass at most
$n/2$ remain; after two, at most $n/4$; after $k$ passes, at most $n/2^k$. The loop
must stop once fewer than one candidate remains, i.e. once $n/2^k < 1$, which is
$2^k > n$, which is $k > \log_2 n$. So the number of passes is at most about
$\log_2 n + 1$.

\begin{complexitybox}
\textbf{Searching a sorted array of size $n$.}
\begin{center}
\begin{tabular}{l c c}
\toprule
 & worst-case comparisons & growth shape \\
\midrule
\textsc{Linear-Search} & $n$ & linear \\
\textsc{Binary-Search} & $\approx \log_2 n + 1$ & logarithmic \\
\bottomrule
\end{tabular}
\end{center}
Binary search needs a \emph{precondition} -- the array must be sorted -- that linear
search does not. That precondition is the price of the speedup, and remembering it is
part of using the algorithm correctly.
\end{complexitybox}

Consider what those numbers mean. For $n = 1{,}000{,}000$, linear search may need a million
comparisons; binary search needs about $\log_2(10^6) \approx 20$. Twenty versus a
million, for the identical answer, on the identical machine. That is not a constant
factor you could buy back with faster hardware -- it is a different \emph{shape} of
growth, and it is the kind of win that good algorithm design hands you over and over.

\begin{trapbox}
Binary search is famous for being easy to get \emph{almost} right and maddening to
get \emph{exactly} right. The classic bugs: writing $lo < hi$ when you meant
$lo \le hi$ (so you miss a one-element slice and report a present value as absent),
or updating $lo \gets mid$ instead of $lo \gets mid + 1$ (so the slice never shrinks
when $mid = lo$, and the loop spins forever). The loop's job is to \emph{strictly
shrink} $hi - lo$ every pass; if a bug lets the slice stand still, termination dies
with it. When you write a loop, always ask: ``what quantity must strictly decrease
to guarantee I stop?''
\end{trapbox}

\begin{walkbox}
Here is a puzzle to carry on a walk. Binary search is gloriously fast \emph{but only
on a sorted array}. Sorting an array, as you will see in the next chapters, costs
roughly $n\log n$ work. So: if you are going to search the array just once, is it
worth sorting it first? What if you will search it a thousand times? Ten thousand?
Try to find the rough break-even point where ``sort once, then binary-search many
times'' beats ``linear-search every time.'' You have just stumbled onto the idea of
\emph{amortizing} a one-time setup cost across many later operations -- a theme that
runs through the whole field.
\end{walkbox}

\section{Algorithms as a technology}

Step back and take in what these two searches taught us. Both were correct. Both ran
on the same hardware. The difference between barely-usable and effortless was not the
machine, not the language, not the programmer's typing speed -- it was the
\emph{idea}: halving instead of scanning. That is what we mean when we say algorithms
are a \emph{technology}, every bit as real as faster processors or bigger memory, and
frequently more decisive.

\begin{aestheticbox}
This is not a toy point about search. The same lesson reappears at the heart of
nearly everything computers do well: the route your map app finds across a continent
in a heartbeat; the way a search engine ranks billions of pages faster than you can
blink; the cryptography that keeps a payment secret; the compression that fits a
movie down a thin wire; the model that recognizes a face. In each, a clever algorithm
turns ``impossible'' into ``instant.'' When a system is too slow, the seasoned
engineer's first question is rarely ``can we buy a faster machine?'' It is ``are we
running the right algorithm?'' More often than newcomers expect, the answer is no --
and fixing it is cheaper, and wins bigger, than any hardware upgrade.
\end{aestheticbox}

\begin{historybox}
The framing of algorithms as a technology, on equal footing with hardware, opens
\textbf{CLRS Chapter~1}. The companion idea -- that we measure an algorithm's quality
by how its resource use \emph{grows} with input size, not by stopwatch readings on
one machine -- is developed in \textbf{CLRS Chapter~2 (``Getting Started'')} and made
rigorous with $O$, $\Theta$, $\Omega$ in \textbf{CLRS Chapter~3
(``Characterizing Running Times'')}. When you reach CLRS, you will find those
chapters move fast; the slow, intuition-first version you are reading now is meant to
make that speed feel comfortable rather than punishing.
\end{historybox}

\section{What to carry forward}

Before the exercises, here is the compressed version -- the handful of ideas worth
keeping in active memory as you read on.

\begin{focusbox}
\textbf{The chapter in seven sentences.}
\begin{enumerate}
  \item An algorithm is a finite, unambiguous, effective sequence of steps from
        input to output -- a recipe with every hidden magic step removed.
  \item A \emph{problem} describes the task for all inputs; an \emph{instance} is one
        concrete input. Correctness, efficiency, and comparison all live at the level
        of the \emph{problem}.
  \item An algorithm is \emph{correct} only if it \emph{halts} with a \emph{right}
        answer on \emph{every} instance; one good test proves nothing, one bad test
        proves it broken.
  \item We argue correctness with a loop invariant via three checks --
        initialization, maintenance, termination -- which is induction in disguise.
  \item We measure efficiency by counting basic operations as a function of input
        size $n$, not by timing seconds on one machine.
  \item What matters is the \emph{growth shape} -- logarithmic, linear, quadratic,
        exponential -- because for large $n$ the shape overwhelms every constant
        factor.
  \item A better algorithm can beat a worse one by a margin that grows with $n$,
        which no faster hardware can match: algorithms are a technology.
\end{enumerate}
\end{focusbox}

Everything ahead -- sorting, data structures, graph algorithms, dynamic programming --
is variations and deepenings of this one melody: specify the problem precisely, design
a method, prove it correct, and analyze how its cost grows. The vocabulary is now in
hand, and the rest of the book puts it to work.

\section*{Exercises}
\addcontentsline{toc}{section}{Exercises}

\begin{exercise}[Wish or algorithm?]
For each ``recipe'' below, decide whether it is a genuine algorithm or a vague wish,
and if it is a wish, name the hidden magic step.
(a) ``To sort the list, put the numbers in increasing order.''
(b) ``To sort the list, repeatedly find the smallest remaining number and move it to
the end of a growing output list, until none remain.''
(c) ``To find a friend in a crowd, look for them.''
(d) ``Start at the left wall; walk forward scanning faces; stop when you see your
friend or reach the right wall.''
\end{exercise}

\begin{exercise}[Problem versus instance]
Write down, in the formal Input/Output style of this chapter, the problem ``find the
\emph{second}-largest element of an array of distinct numbers.'' Then give two
specific instances of size $5$ together with their correct outputs.
\end{exercise}

\begin{exercise}[Trace it]
Trace \textsc{Binary-Search} on $A = \langle 2,5,8,12,16,23,38,56,72,91\rangle$
searching for $v = 12$, and again for the absent value $v = 40$. Produce the full
$(lo, hi, mid, A[mid])$ table for each, and report how many passes each search takes.
\end{exercise}

\begin{exercise}[Off-by-one autopsy]
Someone writes binary search with the loop condition \texttt{while} $lo < hi$
(strict) instead of $lo \le hi$. Find a concrete sorted array and a target $v$
\emph{present} in it for which this buggy version returns \textsc{Nil}. Explain in
one sentence which slice the bug fails to inspect.
\end{exercise}

\begin{exercise}[Counting comparisons]
Exactly how many comparisons $A[i] = v$ does \textsc{Linear-Search} perform, in the
worst case, on an array of size $n$? In the best case? Describe an instance achieving
each.
\end{exercise}

\begin{exercise}[An invariant of your own]
Consider \textsc{Find-Min}, the obvious twin of \textsc{Find-Max} that returns the
\emph{smallest} element. Write its pseudocode, state a loop invariant for it, and
carry out all three checks (initialization, maintenance, termination).
\end{exercise}

\begin{exercise}[Sum of an array]
Write pseudocode for \textsc{Array-Sum}$(A[1\,..\,n])$ returning
$A[1] + A[2] + \cdots + A[n]$. State a loop invariant about the running total, prove
correctness with the three checks, and give the exact number of additions performed
as a function of $n$.
\end{exercise}

\begin{exercise}[Growth-shape intuition]
A quadratic algorithm does $n^2$ steps; a linear one does $100n$ steps -- a hundred
times more work \emph{per element}. For which values of $n$ is the linear algorithm
actually faster? Solve $100n < n^2$ for $n$. What does your answer say about trusting
small-input benchmarks when choosing between algorithms of different growth shapes?
\end{exercise}

\begin{exercise}[Sort-then-search break-even]
Pursue the walking question from the chapter. Suppose sorting an $n$-element array
costs about $6\,n\log_2 n$ basic operations, one linear search costs about $n$, and
one binary search costs about $\log_2 n$. You will perform $k$ searches on the same
array. Write an inequality, in terms of $n$ and $k$, for when ``sort once, then do
$k$ binary searches'' beats ``do $k$ linear searches.'' For $n = 1{,}000{,}000$,
roughly how large must $k$ be? (An estimate is fine.)
\end{exercise}

\begin{exercise}[Halving, generalized -- challenge]
Binary search throws away \emph{half} the candidates per probe. Imagine a
``ternary'' variant that uses two probes per pass to throw away \emph{two thirds} of
the candidates. After $k$ passes, at most how many candidates remain, as a function
of $n$ and $k$? How many passes until fewer than one remains? Compare the pass count
to binary search's $\log_2 n$, and explain why doing more comparisons \emph{per pass}
does not necessarily mean fewer comparisons \emph{overall}.
\end{exercise}

\begin{exercise}[Find the missing number -- challenge]
You are given an array containing $n-1$ distinct numbers, which are some $n-1$ of the
$n$ integers $\{1, 2, \dots, n\}$ in arbitrary order; exactly one integer from that
range is missing. Design an algorithm that finds the missing integer using a number
of basic operations that grows only \emph{linearly} in $n$, \emph{without} sorting
the array. (Hint: what is $1 + 2 + \cdots + n$, and can you compute it with a
formula?) Write the pseudocode and justify why it is correct.
\end{exercise}
